% GENERATED FILE. Edit canonical lessons, then run npm run generate:books.
% canonical-source-hash: fnv1a64:aa12fe7a8374b7f7
% canonical-lessons: GE-C306-gratulieren, GE-C306-glueckwunsch, GE-C306-ueberreichen, GE-C306-beschenken, GE-C306-auspacken

\chapter{Herzlichen Glückwunsch}
\label{ch:ge-herzlichen-gluckwunsch}
\hlchaptermodality{306}

\begin{chapteropening}
\textbf{By the end of this chapter:} \emph{I can talk about congratulations and presents: to congratulate, good wishes, to present, to give presents to and to unwrap.}

Complete the last lesson of chapter 306: I can talk about congratulations and presents: to congratulate, good wishes, to present, to give presents to and to unwrap.
\end{chapteropening}

\section[gratulieren]{gratulieren — to congratulate}

\label{lesson:GE-C306-gratulieren}



\begin{quote}
Before the new one: say the German for the ceremony, then the German for to open officially.
\end{quote}

\subsection*{You'll want to know: gratulieren}
\textbf{gratulieren} — \textquotedblleft{}to congratulate\textquotedblright{}.

\begin{quote}
Alle gratulierten ihr zum Geburtstag.
\end{quote}

\begin{quote}
\emph{Everyone congratulated her on her birthday.}
\end{quote}

With the dative and \emph{zu}: \emph{Ich gratuliere dir zum Abitur}.

\begin{cousinweb}
\emph{gratulieren} is a learned borrowing from Latin \emph{grātulārī} \textquotedblleft{}to wish joy\textquotedblright{}, built on \emph{grātus} \textquotedblleft{}pleasing, thankful\textquotedblright{} (Lewis \& Short); EtymDB lists no parent.
\end{cousinweb}

\begin{grammarlens}[title={What you've built}]
One more word for telling, explaining and describing at B1.
\end{grammarlens}

\subsection*{Your turn}
\begin{itemize}
\raggedright
  \item \emph{Say it:} \emph{gratulieren}
  \item \emph{Say it:} \emph{gratulieren}, once more
  \item \emph{Say it:} a sentence of your own with \emph{gratulieren}
  \item \emph{Recall:} say the German for the ceremony, then the German for to open officially, then say \emph{gratulieren} again
\end{itemize}

\begin{tcolorbox}[breakable,colback=teal!4,colframe=teal!35!black,title={Before you move on}]
What does gratulieren mean? (\textquotedblleft{}To congratulate\textquotedblright{}.) Say it once more.
\end{tcolorbox}
\section[der Glückwunsch]{der Glückwunsch — the good wishes, congratulations}

\label{lesson:GE-C306-glueckwunsch}



\begin{quote}
Before the new one: say the German for to open officially, then the German for to congratulate.
\end{quote}

\subsection*{You'll want to know: der Glückwunsch}
\textbf{der Glückwunsch} — \textquotedblleft{}the good wishes, congratulations\textquotedblright{}.

\begin{quote}
Er schickte ihr einen Glückwunsch zum Geburtstag.
\end{quote}

\begin{quote}
\emph{He sent her birthday wishes.}
\end{quote}

Plural: \emph{die Glückwünsche}. \emph{Herzlichen Glückwunsch!} — congratulations!

\begin{cousinweb}
\emph{Glückwunsch} joins \emph{Glück} \textquotedblleft{}luck\textquotedblright{}, Middle High German \emph{gelücke} (EtymDB), and the inherited \emph{Wunsch} \textquotedblleft{}wish\textquotedblright{}.
\end{cousinweb}

\begin{grammarlens}[title={What you've built}]
One more word for telling, explaining and describing at B1.
\end{grammarlens}

\subsection*{Your turn}
\begin{itemize}
\raggedright
  \item \emph{Say it:} \emph{der Glückwunsch}
  \item \emph{Say it:} \emph{der Glückwunsch}, once more
  \item \emph{Say it:} a sentence of your own with \emph{der Glückwunsch}
  \item \emph{Recall:} say the German for to open officially, then the German for to congratulate, then say \emph{der Glückwunsch} again
\end{itemize}

\begin{tcolorbox}[breakable,colback=teal!4,colframe=teal!35!black,title={Before you move on}]
What does der Glückwunsch mean? (\textquotedblleft{}The good wishes, congratulations\textquotedblright{}.) Say it once more.
\end{tcolorbox}
\section[überreichen]{überreichen — to present, to hand over}

\label{lesson:GE-C306-ueberreichen}



\begin{quote}
Before the new one: say the German for to congratulate, then the German for the good wishes.
\end{quote}

\subsection*{You'll want to know: überreichen}
\textbf{überreichen} — \textquotedblleft{}to present, to hand over\textquotedblright{}.

\begin{quote}
Der Chef überreichte ihm ein Geschenk.
\end{quote}

\begin{quote}
\emph{The boss presented him with a gift.}
\end{quote}

Not separable: \emph{sie hat überreicht}. Formal.

\begin{cousinweb}
\emph{überreichen} joins \emph{über} \textquotedblleft{}over\textquotedblright{} and \emph{reichen} \textquotedblleft{}to hand, to reach\textquotedblright{}, from Old High German \emph{reihhen} (EtymDB), a cousin of English \emph{reach}.
\end{cousinweb}

\begin{grammarlens}[title={What you've built}]
One more word for telling, explaining and describing at B1.
\end{grammarlens}

\subsection*{Your turn}
\begin{itemize}
\raggedright
  \item \emph{Say it:} \emph{überreichen}
  \item \emph{Say it:} \emph{überreichen}, once more
  \item \emph{Say it:} a sentence of your own with \emph{überreichen}
  \item \emph{Recall:} say the German for to congratulate, then the German for the good wishes, then say \emph{überreichen} again
\end{itemize}

\begin{tcolorbox}[breakable,colback=teal!4,colframe=teal!35!black,title={Before you move on}]
What does überreichen mean? (\textquotedblleft{}To present, to hand over\textquotedblright{}.) Say it once more.
\end{tcolorbox}
\section[beschenken]{beschenken — to give presents to}

\label{lesson:GE-C306-beschenken}



\begin{quote}
Before the new one: say the German for the good wishes, then the German for to present.
\end{quote}

\subsection*{You'll want to know: beschenken}
\textbf{beschenken} — \textquotedblleft{}to give presents to\textquotedblright{}.

\begin{quote}
Die Großeltern beschenkten die Kinder reich.
\end{quote}

\begin{quote}
\emph{The grandparents gave the children lots of presents.}
\end{quote}

The person is the object: \emph{jemanden beschenken}.

\begin{cousinweb}
\emph{beschenken} is built on \emph{schenken} \textquotedblleft{}to give as a present\textquotedblright{}, from Old High German \emph{scenken} \textquotedblleft{}to pour in, to give drink to\textquotedblright{} (EtymDB): first a drink poured for a guest, then any gift.
\end{cousinweb}

\begin{grammarlens}[title={What you've built}]
One more word for telling, explaining and describing at B1.
\end{grammarlens}

\subsection*{Your turn}
\begin{itemize}
\raggedright
  \item \emph{Say it:} \emph{beschenken}
  \item \emph{Say it:} \emph{beschenken}, once more
  \item \emph{Say it:} a sentence of your own with \emph{beschenken}
  \item \emph{Recall:} say the German for the good wishes, then the German for to present, then say \emph{beschenken} again
\end{itemize}

\begin{tcolorbox}[breakable,colback=teal!4,colframe=teal!35!black,title={Before you move on}]
What does beschenken mean? (\textquotedblleft{}To give presents to\textquotedblright{}.) Say it once more.
\end{tcolorbox}
\section[auspacken]{auspacken — to unwrap, to unpack}

\label{lesson:GE-C306-auspacken}



\begin{quote}
Before the new one: say the German for to present, then the German for to give presents to.
\end{quote}

\subsection*{You'll want to know: auspacken}
\textbf{auspacken} — \textquotedblleft{}to unwrap, to unpack\textquotedblright{}.

\begin{quote}
Das Kind packte die Geschenke sofort aus.
\end{quote}

\begin{quote}
\emph{The child unwrapped the presents at once.}
\end{quote}

Separable. The opposite of \emph{einpacken}, already known.

\begin{cousinweb}
\emph{auspacken} joins \emph{aus} \textquotedblleft{}out\textquotedblright{} and \emph{packen} \textquotedblleft{}to pack\textquotedblright{}, which was borrowed from Middle Low German \emph{packen} (EtymDB).
\end{cousinweb}

\begin{grammarlens}[title={What you've built}]
One more word for telling, explaining and describing at B1.
\end{grammarlens}

\subsection*{Your turn}
\begin{itemize}
\raggedright
  \item \emph{Say it:} \emph{auspacken}
  \item \emph{Say it:} \emph{auspacken}, once more
  \item \emph{Say it:} a sentence of your own with \emph{auspacken}
  \item \emph{Recall:} say the German for to present, then the German for to give presents to, then say \emph{auspacken} again
\end{itemize}

\begin{tcolorbox}[breakable,colback=teal!4,colframe=teal!35!black,title={Before you move on}]
What does auspacken mean? (\textquotedblleft{}To unwrap, to unpack\textquotedblright{}.) Say it once more.
\end{tcolorbox}
